\section{Impulso y cantidad de movimiento}

\subsection{Cantidad de movimiento lineal}
\begin{equation}
    \vec{p} = m \vec{v}
\end{equation}
Unidad: \(\text{kg} \cdot \text{m/s}\)

\subsection{Impulso}
\begin{equation}
    \vec{I} = \vec{F} \cdot \Delta t = \Delta \vec{p}
\end{equation}

\subsection{Teorema del impulso y cantidad de movimiento}
\begin{equation}
    \vec{F}_{\text{prom}} \Delta t = m \vec{v}_f - m \vec{v}_i
\end{equation}

\section{Conservación de la cantidad de movimiento}
Para un sistema aislado (sin fuerzas externas netas):
\begin{equation}
    \vec{p}_{\text{total, inicial}} = \vec{p}_{\text{total, final}}
\end{equation}

\section{Tipos de colisiones}

\subsection{Colisiones elásticas}
Se conservan la cantidad de movimiento y la energía cinética.
\begin{align}
    m_1 v_{1i} + m_2 v_{2i} &= m_1 v_{1f} + m_2 v_{2f} \\
    \frac{1}{2} m_1 v_{1i}^2 + \frac{1}{2} m_2 v_{2i}^2 &= \frac{1}{2} m_1 v_{1f}^2 + \frac{1}{2} m_2 v_{2f}^2
\end{align}

\subsection{Colisiones inelásticas}
Se conserva la cantidad de movimiento pero \textbf{no} la energía cinética (se transforma en calor, sonido, etc.).
\begin{equation}
    m_1 v_{1i} + m_2 v_{2i} = (m_1 + m_2) v_f \quad (\text{perfectamente inelástica})
\end{equation}

\resumebox{ejercicio}{Fuerza magnética del protón}{\footnotesize
    Una bala de \(10 \, \text{g}\) se dispara contra un bloque de \(2 \, \text{kg}\) en reposo. Después del impacto, la bala queda incrustada y el sistema se mueve a \(0.5 \, \text{m/s}\). Calcula la velocidad inicial de la bala.
}